%==============================================================================
%  sections_v3/theorem_section.tex -- the results section of draft_v3.
%
%  SOURCE OF TRUTH: research/model_v4/MODEL_CARD.md, version stamp 2026-08-28
%  (re-review audit repairs, commit 59c0dfc), section 6
%  result ledger.  The eight
%  results appear in the ledger's own order and at its grain: D1, L1, L2, L3,
%  L4, P1, T1, C1.  Every hypothesis clause of every card row travels with the
%  statement that carries it; the crosswalk in the ticket 08 scratchpad walks
%  them one by one.  Honesty labels are transcribed verbatim from the ledger,
%  conditionality included.
%
%  Proofs live in sections_v3/proofs_section.tex; no proof environment appears
%  in this file.
%==============================================================================

\section{Results}
\label{sec:results}

Eight results follow, in the order the model builds them: the partition and its accounting
identities, then what liquidity does inside each cell, then existence, then the attenuation theorem,
then what survives the move from fixed policies to equilibrium. Each statement carries the honesty
label of the record it comes from, and each carries that label's conditionality with it. The
conditionality is not decoration; where a hypothesis has been measured and found to fail at the
calibration the accompanying implementation runs, the statement says so beside itself rather than in
an appendix.

\subsection{The partition and its accounting identities}
\label{sec:partition-results}

\begin{lemma}[D1: the disclosure partition and the price-path decomposition]
\label{lem:D1}
Assume Assumptions~\ref{asm:A1}, \ref{asm:A2p}, \ref{asm:A4} and \ref{asm:A5}; the primitive table
restrictions of Assumption~\ref{asm:TR}; the cutoff selection map of Definition~\ref{def:cpbe}\,(i); and the two price
conventions of Definition~\ref{def:prices}, namely $P_{-1}^P = \Eop[Y]$ and
$P_{\ND} = P_{f^-}^P$. Assume in addition the two clauses now carried as primitives: Borel
regularity of $s \mapsto B_j(s,d)$ for \emph{every} plan including Exit, clause (TR-ii), which is
needed only for part~(c) below because pooled pricing integrates over every type; and a content for
the control-node information set $\mathcal I_H$, supplied by Definition~\ref{def:prices}. Then:
\begin{enumerate}[label=(\alph*),leftmargin=2.4em,itemsep=3pt]
\item the disclosure indicator $D = \ind\{a = 1,\ c(\tau) + T \le H\}$ is measurable, and it maps
  every control-node history into exactly one cell;
\item for every Voice plan, $f_j \le H$ if and only if $B_j(s, H-T) \ge \tau$;
\item each flagged history yields $B^F$, $R_d$, $R$ and $J$, and these satisfy the exact identity
  \begin{equation}
  \label{eq:runup-jump}
  P^F - P_{c^-}^P \;=\; R + J.
  \end{equation}
\end{enumerate}
\end{lemma}
\resultstatus{PROVED}

The proof is in Section~\ref{sec:proofs}. Part~(b) is the \emph{clock equivalence} that
Lemma~\ref{lem:L4} and Theorem~\ref{thm:T1} both consume: it converts a statement about the filing
date into a statement about the stake path $T$ days before the horizon, which is what makes window
comparisons tractable.

\begin{lemma}[L1: two-cell decomposition of the engagement premium]
\label{lem:L1}
Assume Lemma~\ref{lem:D1}; the definitions of Definitions~\ref{def:prices} and \ref{def:premium};
Assumption~\ref{asm:A5}, which pins \emph{the} version of $\Eop[Y \mid \mathcal I]$;
Assumption~\ref{asm:A2p} together with clause (TR-i), under which $\Delta_m$ is finite; and
Assumption~\ref{asm:A1}, which puts every object on one probability space. Then whenever
$0 < \Omega < 1$,
\begin{equation}
\label{eq:L1}
\Dact \;=\; \Omega\,M_F \;+\; (1-\Omega)\,M_P .
\end{equation}
At $\Omega = 1$ the identity degenerates to $\Dact = M_F$, and at $\Omega = 0$ to $\Dact = M_P$; in
each degenerate case the average over the null cell is \emph{undefined rather than imputed}. That
last clause is a non-identification statement, and it is proved rather than asserted.
\end{lemma}
\resultstatus{PROVED}

The proof is in Section~\ref{sec:proofs}.

\subsection{Liquidity in the flagged cell}
\label{sec:flagged-results}

\begin{lemma}[L2: $\kappa$-invariance of the flagged cell]
\label{lem:L2}
At fixed cutoff \emph{and} execution policies, assume Assumption~\ref{asm:A1};
Assumption~\ref{asm:A2p} \emph{together with} the primitive table restrictions of
Assumption~\ref{asm:TR}; Assumption~\ref{asm:A4}; Assumption~\ref{asm:A5};
Assumption~\ref{asm:A7p} in its on-path injective form, consumed almost surely on the flagged
set;\footnote{``Almost surely'' is the permissive reading, and it is the only coherent one when
$B^F$ is continuum-valued, since then no individual flagged tuple has positive probability.}
Lemma~\ref{lem:D1}; the no-feedback timing of Assumption~\ref{asm:nofeedback}; and $\Omega > 0$.
Assume also an explicit bidder-entry rule --- either \eqref{eq:entry} or any rule with the two
properties named in the proof --- which is carried as bookkeeping. Then the flagged tuple
$\Sfl = (B^F, Q^F, a{=}1)$ makes the pre-filing pooled history conditionally independent of
$(v,s,\xi)$ on the flagged set, and consequently the flagged posterior, the flagged price, the entry
probability and $M_F$ are all invariant to $\kappa$.
\end{lemma}
\resultstatus{PROVED}

The proof is in Section~\ref{sec:proofs}. Assumption~\ref{asm:A7} in its weak identification wording
is not sufficient here: it permits two pairs $(j,s)$ with different pooled paths, which is this
lemma's first failure case. The form consumed is Assumption~\ref{asm:A7p}, named as such.

\subsection{Liquidity in the pooled cell}
\label{sec:pooled-results}

\begin{lemma}[L3: interior motion of the pooled premium under the chord restriction]
\label{lem:L3}
Assume $\Atau$ (Assumption~\ref{asm:Atau}), including both of its clauses: ($\tau$-i),
kernel-through-posterior, and ($\tau$-ii), $\kappa$-free support \emph{and} $\kappa$-free $\bar\pi$.
Assume in addition: $h(0) = 0$; $\kappa$-free pooled mass and $\kappa$-free pooled engagement moment
at fixed policies; Lemma~\ref{lem:D1} by statement; and regularity stated minimally --- $h$
continuous on $[0,\bar\pi]$ and twice differentiable on the open interval $(0,\bar\pi)$, Darboux's
theorem doing the rest, with no continuity of $h''$ required. For the small-$\bar\pi$ corollary
(part~(c)) assume also a second-order Peano expansion of $h$ at $0+$ and one and the same kernel
along the shrinking family; and for the seam at which Lemma~\ref{lem:L4} consumes this lemma, assume
$|\Akap|$ bounded \emph{uniformly in $\bar\pi$} along the limit. Then:
\begin{enumerate}[label=(\alph*),leftmargin=2.4em,itemsep=3pt]
\item $\partial_\kappa \Eop_\kappa[h] = \Akap\,C_h(\bar\pi)$, exactly;
\item $C_h(\bar\pi) = \tfrac14\bar\pi^2 h''(\zeta)$ for some $\zeta \in (0,\bar\pi)$ --- an identity,
  not an approximation;
\item $C_h = \tfrac14 h''(0)\,\bar\pi^2 + o(\bar\pi^2)$, so the interior motion vanishes at rate
  $\bar\pi^2$ as $\bar\pi \downarrow 0$.
\end{enumerate}
The statement is an ``if'' and never an ``iff'': $\Akap = 0$ also kills the interior motion.
\end{lemma}
\resultstatus{PROVED under $\Atau$}

The proof is in Section~\ref{sec:proofs}. The conditionality is the substance of the result, not a
caveat on it. Whether the two-round pooled cell of Section~\ref{sec:timing} satisfies the support
condition of $\Atau$ is \textbf{open}: the discussion after Assumption~\ref{asm:Atau} shows that the
support condition is the assumption's entire remaining bite, and exhibits a one-round market that
satisfies it and a no-disclosure structure that does not; the weakest sufficient conditions for the
two-round case are named in the closing steps of the proof in
Section~\ref{sec:proofs}. Remark~\ref{rem:AtauRecord} records that at the implemented calibration the
support condition fails at all $180$ non-degenerate nodes; at that calibration this lemma therefore
says nothing about the implemented pooled cell, and its label is untouched by that fact.

\begin{lemma}[L4: threshold tightening --- nesting, composition, and the pooled sensitivity]
\label{lem:L4}
At fixed policies, for $b_0 < \tau' < \tau$ at a common window $T$ and a common $\kappa$, with
$\Omega(\tau',T) < 1$:
\begin{enumerate}[label=(leg \arabic*),leftmargin=3.4em,itemsep=3pt]
\item \emph{(unconditional)} $\mathcal C_F(\tau,T) \subseteq \mathcal C_F(\tau',T)$, with every newly
  flagged history generated by a Voice plan; hence $\Omega(\tau',T) \ge \Omega(\tau,T)$.
\item \emph{(unconditional)} The pooled engagement \emph{share} falls,
  $\bar\pi_{\mathrm{pr}}(\tau') \le \bar\pi_{\mathrm{pr}}(\tau)$, with an exact identity for the
  difference between the two shares.
\item \emph{(under $\Abr$)} $\mathcal S_P(\tau',T) \le \mathcal S_P(\tau,T)$, with equality whenever
  $C_h(\bar\pi(\tau)) = 0$.
\end{enumerate}
Legs 1 and 2 need only: the clock equivalence of Lemma~\ref{lem:D1}\,(b); the no-feedback timing of
Assumption~\ref{asm:nofeedback}; fixed policies; $b_0 < \tau' < \tau$ imposed at \emph{both}
thresholds; Assumptions~\ref{asm:A1} and \ref{asm:A4}; clause (TR-iii)'s implication
$D = 1 \Rightarrow a = 1$; and $\Omega(\tau') < 1$. Nestedness is a \emph{conclusion} of leg 1, not a
hypothesis of it. Leg 3 additionally needs: Lemma~\ref{lem:L3} by statement; the maintained
\emph{magnitude} monotonicity of $|C_h|$ in $\bar\pi$ from Assumption~\ref{asm:Atau}, the sign half
$C_h \le 0$ being unused at this leg; and clauses (br-i)--(br-v) of Assumption~\ref{asm:Abr}.
\end{lemma}
\resultstatus{PROVED (legs 1--2 outright; leg 3 under $\Abr$)}

The proof is in Section~\ref{sec:proofs}. Leg 3 inherits the conditionality of
Assumption~\ref{asm:Abr}, and through clause (br-i) it inherits the representation \eqref{eq:atau} at
\emph{both} compared thresholds; Remark~\ref{rem:AtauRecord} therefore bears on leg 3 as directly as
it bears on Lemma~\ref{lem:L3}. At the implemented calibration leg 3's antecedent is not satisfied,
and at that calibration leg 3 says nothing about the implemented pooled cell. Legs 1 and 2 are
unaffected: they consume no chord restriction. Clause (br-iii) is the clause with the least
justification behind it and is the one to attack first.

\subsection{Existence}
\label{sec:existence-results}

The existence result is stated as a proposition rather than as a theorem: the paper has one theorem,
Theorem~\ref{thm:T1}, and existence is the standing-ground for it rather than the claim the paper is
making. Its hypothesis set is long, and every clause of it is load-bearing --- the 2026-08
demotion and restoration of this result turned on exactly which form of one hypothesis the argument
consumed --- so the hypotheses are enumerated rather than gathered into a sentence.

\begin{proposition}[P1: existence of a cutoff perfect Bayesian equilibrium]
\label{prop:P1}
Assume:
\begin{enumerate}[label=(P-\arabic*),leftmargin=3.2em,itemsep=3pt]
\item Assumption~\ref{asm:A1} (independent primitives);
\item Assumption~\ref{asm:A2p} (finite model, integrable objective);
\item Assumption~\ref{asm:A3} (ordered plans, single crossing);
\item Assumption~\ref{asm:A4} (legal-clock discipline);
\item Assumption~\ref{asm:A6} (compact outer self-map), read as in the second trailing paragraph
  below;
\item Assumption~\ref{asm:A7J}, \emph{joint tuple injectivity}: the map
  $(j,s) \mapsto (B_j^F, Q_j^F, a_j)$ is injective on the whole flagged-pair set
  $\{(j,s) : D_j = 1\}$, \emph{including pairs no cutoff vector selects}. This is strictly stronger
  than the on-path form, Assumption~\ref{asm:A7p}, and it is the form the argument consumes where it
  pins \emph{off-path} flagged beliefs;\footnote{The distinction is not cosmetic. Before 2026-08-25
  the recorded statement of this result carried the on-path form while the proof consumed the joint
  form, so the two verification passes on file covered two different statements; that mismatch is
  what the 2026-08-23 demotion turned on, and the statement here is the amended one, re-verified by
  a fresh adversarial proof-read and an independent statements-only re-derivation, both 2026-08-25
  and both by agents who did not write the proof. The form is named every time either appears.}
\item Lemma~\ref{lem:D1} by statement, \emph{with its own hypotheses travelling};
\item the no-feedback timing of Assumption~\ref{asm:nofeedback}, read together with the
  flag-termination clause of Assumption~\ref{asm:flagterm};
\item the \emph{definitional} round-2 action-set hypothesis: the flagged-round action set \emph{is}
  the plan-generated set $\{Q^F_{j'}(s)\}$ taken over menu elements agreeing with $j$ on everything
  already played. This is not a closure condition; the closure form is jointly unsatisfiable with
  finiteness, by cardinality;
\item \emph{continuation-cost equivalence on that same set}: menu elements sharing $j$'s pooled path
  up to $f_j(s)$ and carrying $a_{j'} = a_j$ have the same engagement cost, $C_{j'}(s) = C_j(s)$.
  This holds trivially on any single-Voice menu, where the set is a singleton, and it is live only on
  menus with two or more Voice plans sharing a pooled path. What it buys, under the
  plan-completion reading of the timing convention in Remark~\ref{rem:Ctiming}, is this: on that set
  the flagged price does not move and the flagged order cancels, so the engagement cost is the only
  thing that can differ between staying and deviating, and at a flagged pair the cutoff vector does
  not select there is no date-$0$ optimality to fall back on --- so without this clause the deviator
  takes the class member with the smallest cost and requirement (ii) of Definition~\ref{def:cpbe}
  fails at that node. Under the sunk reading the continuation is constant on the deviation set with
  no clause at all and this hypothesis is not consumed; it is listed because the statement does not
  commit to a reading, and it is what makes the conclusion hold under both;
\item the restriction $m_0 \ge 0$ of \eqref{eq:m0};
\item the blockholder-objective definition of Definition~\ref{def:U}, whose $-a_j C_j(s)$ term is
  the display \eqref{eq:Uj}, together with the timing convention for $C_j(s)$ in
  Remark~\ref{rem:Ctiming} --- the engagement cost may be booked on completing the plan or as sunk
  once the filing has landed, the two give the same round-2 comparison on the round-2 deviation set,
  and the result does not depend on the choice;
\item the primitive table restrictions \ref{asm:TR} that the argument consumes, in particular: the
  terminal payoff \eqref{eq:Y} with the price convention $P(\mathcal I) = \Eop[Y \mid \mathcal I]$
  and the entry rule \eqref{eq:entry}, both from clause (TR-iv); clause (TR-ii)'s Borel regularity
  for \emph{every} plan including Exit, needed here \emph{directly} and not by way of
  Lemma~\ref{lem:D1}, whose conclusion is measurability of $D$ and the cell map; clause (TR-iii)'s
  implication $D = 1 \Rightarrow a = 1$, the definitions of $c$, $f$, $B^F$, $Q^F$ and $b^*$, and
  $\partial_s B_j \ge 0$ for Voice; and clause (TR-i)'s distributional forms with $\Delta_m > 0$.
\end{enumerate}
Then \emph{at every $\kappa \in [0,1]$} a cutoff perfect Bayesian equilibrium over complete
contingent plans exists: there is $k^\star \in \Theta$ with $k^\star = \Tmap(k^\star;\vartheta)$,
prices at their inner fixed points and beliefs Bayes-consistent on path, together with the following.
\begin{enumerate}[label=(\roman*),leftmargin=2.4em,itemsep=3pt]
\item Off-path beliefs are limits of \emph{one} full-support perturbation family over plans, fixed
  once and used to define the price system at every $k \in \Theta$ and not only at $k^\star$ ---
  because the deviation payoffs that define $\Tmap$ read off-path pooled histories --- at every
  pooled history reachable \emph{with positive probability} under some plan profile.
\item The boundary values of $\kappa$ are handled by \emph{extension}, not by restriction. At
  $\kappa \in \{0,1\}$ the noise support degenerates to $\{0\}$ and to $\{\pm\bar z\}$ respectively;
  a pooled history needing a mark outside the surviving support is null under \emph{every} profile,
  so it is off nature's path rather than off the players', carries no requirement under
  Definition~\ref{def:cpbe}\,(vi), and is read by no step. No cut to $\kappa \in [0,1)$ is taken, and
  the earlier claim of a belief at \emph{every} pooled history --- false at $\kappa = 1$ --- is
  withdrawn.
\item Flagged-tuple beliefs are supplied by Assumption~\ref{asm:A7J} at every tuple in the image of
  the flagged-pair map $(j,s) \mapsto (B_j^F,Q_j^F,a_j)$, on path and off, since the image includes
  tuples generated by pairs the cutoff vector does not select. The belief is the point mass at the
  unique generating pair. That point mass is a \emph{version} of the conditional law at every image
  tuple --- the signal is continuous, so a version is what a conditional law is, and any a.e.-equal
  version serves requirements (iii) and (vi) of Definition~\ref{def:cpbe} equally --- and it is the
  version this equilibrium selects. No tuple outside that image arises, because the round-2
  action-set hypothesis (P-9) leaves no off-menu order to produce one.
\item The bidder follows the entry rule \eqref{eq:entry}.
\item There is a sequentially optimal flagged component at \emph{every} flagged pair $(j,s)$,
  whether or not the cutoff vector selects it. The flagged price is invariant across the round-2
  deviation set, since Assumption~\ref{asm:A7J} pins the belief at the same $s$ and at $\pi = 1$, so
  the flagged order cancels out of the continuation and hypothesis (P-10) makes what remains
  constant.
\end{enumerate}
Two readings belong to the statement rather than to the proof. First, \emph{Assumption~\ref{asm:A5}
is not assumed here}: its existence and uniqueness content is derived from $m_0 \ge 0$, its
continuity content \emph{in the belief summaries} $(\hat v,\pi)$ from the same scalar reduction, and
its measurable-selection content from Assumption~\ref{asm:A7J} together with clause (TR-ii)'s Borel
regularity. What is \emph{not} derived is the cutoff clause of Assumption~\ref{asm:A5}: continuity
of the composed pooled price family in the cutoff vector $k$, which runs through the conditioning
pair $(\hat v,\pi)$ rather than through the pricing map, as Step~7 of the proof says where it says
what is left over. That continuity enters only through the reading of Assumption~\ref{asm:A6} given
next, and Remark~\ref{rem:A6record} records it measured to fail at the implemented calibration.
Second,
\emph{Assumption~\ref{asm:A6} is read} as asserting that $\Tmap$ --- under a named
tie-break-and-corner selection, without which a correspondence cannot be called continuous --- is a
well-defined, single-valued, continuous self-map of $\Theta$, with $\Theta$ nonempty as in
Definition~\ref{def:theta}.

Finally, \emph{at any such equilibrium at which Assumption~\ref{asm:A8} holds}, both cells carry
strictly positive probability and both are on path; and for the restatement of
Assumption~\ref{asm:A8} as a single signal threshold, add H-ord --- Voice stake monotonicity across
plans --- and the upper-set engagement-flag hypothesis. \emph{Uniqueness is not claimed.}
\end{proposition}
\resultstatus{PROVED}

The proof is in Section~\ref{sec:proofs}.

\begin{remark}[Numerical record for Proposition~\ref{prop:P1}]
\label{rem:P1record}
The label above rests on the proof and on the two 2026-08-25 verification passes, not on a grid, and
the grid evidence is stated separately for that reason. Four nodes of the equilibrium sweep ---
$\kappa \in \{0.15,\,0.85\}$ crossed with $(\tau,T) \in \{(0.05,\,5),\,(0.075,\,1)\}$ --- remain
STILL UNRESOLVED after $30$ seeds each: the best payoff-scale residual runs
$3.1\times10^{-4}$ to $1.5\times10^{-3}$ against a $10^{-9}$ criterion, while the best cutoff-scale
residual runs $10^{-14}$ to $10^{-11}$. The A3 and A6 proxies pass at every achieving seed, which,
as Remark~\ref{rem:A3record} explains, is not evidence that those hypotheses hold: the proxies are
local screens and are silent on argmax monotonicity over $s$ and on continuity of $\Tmap$ in $k$.
This evidence is UNCHECKED: existence at those four nodes is neither claimed nor denied by
it. Record: \texttt{quality\_reports/fixes/t2\_p1\_fournode\_recheck.json}.

Two hypotheses of this proposition carry their own measured failures at the implemented calibration.
Assumption~\ref{asm:A6}'s continuity clause fails there for the $\Theta$ this argument constructs
(Remark~\ref{rem:A6record}), and Assumption~\ref{asm:A3} fails there at two independently found loci
(Remark~\ref{rem:A3record}). Neither moves a label, and none is licensed to: both are hypotheses.
The proposition stays PROVED as a conditional, in the same pattern as $\Atau$ --- what is on
record is that its antecedent, read with the $\Theta$ the proof constructs, is not satisfied by the
implemented calibration, so at that calibration the proposition asserts nothing about the implemented
cell. Nonexistence is neither claimed nor shown anywhere: $23$ of $27$ sweep nodes converge, and two
fixed points survive even at the worst probed node.
\end{remark}

\subsection{Disclosure attenuation}
\label{sec:attenuation}

\begin{theorem}[T1: disclosure attenuation at fixed policies]
\label{thm:T1}
At fixed plan and cutoff policies, with $0 < \Omega < 1$ and $\mathcal S_P > 0$, and under the
hypotheses (T-1)--(T-15) listed below:
\begin{enumerate}[label=(\Alph*),leftmargin=2.4em,itemsep=3pt]
\item \emph{Factorisation.} $\mathcal S = (1-\Omega)\,\mathcal S_P$, exactly; and the same
  factorisation holds for the total-variation aggregate of $\Dact$ over any $\kappa$-grid, with no
  differentiability required.
\item \emph{Threshold margin.} Threshold tightening attenuates:
  \begin{equation}
  \label{eq:T1B}
  \frac{\mathcal S(\tau')}{\mathcal S(\tau)} \;=\; W_\tau\,C_\tau \;\le\; 1,
  \end{equation}
  because \emph{both} ratios lie in $[0,1]$. No dominance condition is needed.
\item \emph{Window margin.} Window tightening attenuates \emph{if and only if} $W_T C_T \le 1$, where
  $W_T \le 1$ is proved --- from the clock equivalence of Lemma~\ref{lem:D1}\,(b) and the monotone
  Voice stake path --- and $C_T$ is \emph{unsigned}. The equivalent form
  \begin{equation}
  \label{eq:T1C}
  \frac{\partial_{r_T}\mathcal S_P}{\mathcal S_P} \;\le\; \frac{\Omega_{r_T}}{1-\Omega}
  \end{equation}
  holds \emph{on average along the tightening path}, integrated over $[-T,-T']$, and exactly in the
  infinitesimal limit; read pointwise, \eqref{eq:T1C} is false.
\end{enumerate}
The hypotheses are: (T-1) fixed policies; (T-2) Assumption~\ref{asm:A8} at each compared policy;
(T-3) $\mathcal S_P > 0$; (T-4) Lemma~\ref{lem:L1}; (T-5) Lemma~\ref{lem:L2}, with its own
hypotheses travelling; (T-6) Lemma~\ref{lem:D1}; (T-7) PE-$\Omega$, that is
$\partial_\kappa\Omega = 0$ at fixed policies --- derivable rather than assumed, and it is exactly
what fails in general equilibrium, which is the term Proposition~\ref{prop:C1} bounds; (T-8)
$\kappa$-differentiability of $M_P$, which no standing hypothesis of Section~\ref{sec:hypotheses}
supplies and which is therefore carried in the proof; (T-9) Assumption~\ref{asm:Atau} at
\emph{both} compared policies, with $\bar\pi$ read as the upper support point; (T-10)
Lemma~\ref{lem:L3}; (T-11) clauses (br-i)--(br-v) of Assumption~\ref{asm:Abr} at the threshold pair;
(T-12) Lemma~\ref{lem:L4}; (T-13) the no-feedback timing of Assumption~\ref{asm:nofeedback}; (T-14) a
smooth window interpolation, for the local form \eqref{eq:T1C} only; and (T-15) threshold-side
smoothness, which has been confirmed non-load-bearing. \emph{No unconditional window sign is
claimed.}
\end{theorem}
\resultstatus{PROVED at fixed policies}

The proof is in Section~\ref{sec:proofs}. Part~(A) is the weight effect in its exact form: the
liquidity sensitivity of the expected engagement-related premium is the pooled cell's sensitivity,
scaled by the pooled cell's weight. Part~(B) is the threshold-margin result and needs no dominance
condition, because the weight ratio and the composition ratio both lie in $[0,1]$. Part~(C) is the
window-margin result, and it is an ``if and only if'' rather than a sign: the composition ratio $C_T$
is unsigned, so the model does not deliver window-margin attenuation as a theorem.

\begin{remark}[Numerical record for Theorem~\ref{thm:T1}]
\label{rem:T1record}
Part~(B) rests on Assumption~\ref{asm:Atau} at both compared policies (T-9), on Lemma~\ref{lem:L3}
(T-10), and on Assumption~\ref{asm:Abr} (T-11). Remark~\ref{rem:AtauRecord} records that the support
condition of $\Atau$ fails at all $180$ non-degenerate nodes of the implemented calibration, and
clause (br-i) carries that representation at both thresholds, so the record bears on (T-9) and (T-11)
alike. Part~(B) stays PROVED as a conditional with its proof untouched; at this calibration
it says nothing about the implemented pooled cell. No label moves.

For Part~(C) the genuine window-margin record is block 4 of the executed T1 check: $W_T C_T < 1$ at
every checked node at this calibration, with an $H = 10$ corner caveat recorded in the handoff.
That is NUMERICAL node evidence at one calibration and not a sign theorem, which Part~(C)
does not claim. One distinction is worth stating because it is easy to lose: the O-1 comparison,
whose ratios $1.06397$, $1.18373$, $1.13631$ and $0.37798$ at
$\Omega = 0.037252$, $0.128950$, $0.285804$ and $0.50$ are sometimes read as a window test, is a
\emph{disclosure-regime} comparison at a fixed filing window in the static repository model. Those
ratios are regime-comparison composition outcomes; they are not $W_T C_T$ and they measure no window
pair. The analogy is useful only because it shows that a composition factor can exceed one, which is
what motivates the genuine window-margin ``if and only if'' above; the O-1 cut at
$\Omega^\star \approx 0.343$ is a disclosure-regime boundary, not a window boundary.
\end{remark}

\subsection{From fixed policies to equilibrium}
\label{sec:ge-results}

Theorem~\ref{thm:T1} holds at fixed policies, and hypothesis (T-7) is exactly what fails once the
cutoff vector is allowed to move with $\kappa$. The next result asks what survives that move. It
carries the region as a \emph{named hypothesis} rather than exhibiting one.

\begin{proposition}[C1: the dominance-and-contraction implication in general equilibrium]
\label{prop:C1}
Assume:
\begin{enumerate}[label=(C-\arabic*),leftmargin=3.2em,itemsep=3pt]
\item the along-the-path contraction of Assumption~\ref{asm:AGE}, $L_{\mathcal R} < 1$, in
  \emph{one fixed norm convention} --- an induced operator norm with its dual pairings; a mismatched
  pairing silently voids the implication;
\item $\mathcal R_r$ relatively open in \emph{both} coordinates, with $\kappa \notin \{0,1\}$;
\item an interior single-branch equilibrium;
\item twice continuous differentiability of $\Dact$ in $(k,\kappa,r)$;
\item a non-vanishing \emph{equilibrium} liquidity derivative --- the equilibrium sensitivity
  $\mathcal S^{GE}$, which is distinguished from the fixed-policy sensitivity $\mathcal S$ of
  Definition~\ref{def:premium};
\item strict dominance $g_r^{PE} > \mathcal B_r^{GE}$ on the region, with $g_r^{PE}$ as in
  \eqref{eq:gPE} and $\mathcal B_r^{GE}$ as in \eqref{eq:BGE};
\item the \emph{threshold} margin $r_\tau$ only: the window coordinate is an integer, and nothing
  local is claimed there.
\end{enumerate}
Then the fixed-policy attenuation sign survives in equilibrium on the region:
\begin{equation}
\label{eq:C1}
\partial_r \mathcal S^{GE} \;\le\; -\eta_r \;<\; 0 .
\end{equation}
The sign-coherence hypothesis --- that the fixed-policy attenuation sign agrees with the equilibrium
sign --- is confirmed unused in \eqref{eq:C1}: its work is one of naming, fixing when that conclusion
may be called a survival of the fixed-policy sign. The sign-constancy clause of
Assumption~\ref{asm:AGE} is likewise unused.
\end{proposition}
\resultstatus{PROVED (dominance-and-contraction implication; region-as-hypothesis) plus NUMERICAL
node evidence}

The proof is in Section~\ref{sec:proofs}.

\begin{remark}[Numerical record for Proposition~\ref{prop:C1}]
\label{rem:C1record}
Eighteen of eighty grid nodes are pointwise \emph{dominance-and-contraction nodes}. The largest
contiguous block is $T = 5$ at $\tau$-percentiles $\{50,70,90\}$ with
$\kappa \in \{0.65,\,0.75,\,0.85\}$; the slack $\eta_r$ has minimum $0.0595$ and median $0.3467$, and
$L_{\mathcal R} \in [0.264,\,0.501]$ everywhere on the grid. These come from the executed committed
check, independently re-run on 2026-08-22, on which all values reproduce. What these nodes verify is
narrow and worth stating exactly: they verify the two pointwise inequalities $L_{\mathcal R} < 1$ and
$\eta_r > 0$ together with supporting diagnostics. They do \emph{not} verify the full antecedent
(C-1)--(C-7) above, and they do \emph{not} exhibit a named nonempty region. A dominance-and-contraction
node is not a fifth honesty label. Record:
\texttt{quality\_reports/fixes/t2\_c1\_region\_check.py} and \texttt{t2\_c1\_region\_check.json},
with the independent re-run verify note beside them. The re-derivation on file adds the bound
$\mathcal B_r^{GE} = O\bigl((1-L_{\mathcal R})^{-3}\bigr)$, so the dominance-and-contraction nodes
are cubically bounded away from $L_{\mathcal R} = 1$.

A named-region promotion is a separate step and is not claimed. The earlier aspiration --- this
result PROVED on a named nonempty region and NUMERICAL off it --- is retired as
structurally undeliverable as worded. What is deliverable is the three objects above: the
implication PROVED with the region as a hypothesis, the dominance-and-contraction nodes
NUMERICAL, and a named-region promotion left open, for which the $\varepsilon$-ball plus
integral-control pattern of the frozen manuscript's D8 is the template.
\end{remark}

\subsection{What is not claimed}
\label{sec:not-claimed}

The results above do not claim: a global window-margin attenuation sign; $\kappa$-invariance of the
filing-day jump $J$; equilibrium uniqueness; a nonempty general-equilibrium region as a theorem;
endogenous filing before the deadline; noisy or partially revealing flagged-round trading;
continuous-time execution; welfare or optimal rule design; that the frozen manuscript's hump result
survives; that the prior calibration at $\Omega \approx 0.037$ is economically meaningful; or any
empirical value for the disclosed share of engagements $\omega_a$.

Three questions remain open, in whole or in part, and are recorded as such. The third was answered in
substance on 2026-08-27; what survives of it is scoped below rather than left standing.
\begin{enumerate}[label=(\arabic*),leftmargin=2.4em,itemsep=3pt]
\item \emph{Whether the two-round pooled cell of Section~\ref{sec:timing} satisfies $\Atau$.}
  Lemma~\ref{lem:L3} shows that the entire remaining bite of the restriction is its support
  condition, exhibits a one-round market that satisfies it and a no-disclosure structure that does
  not, and names the weakest sufficient conditions for the two-round case. Lemma~\ref{lem:L3},
  leg 3 of Lemma~\ref{lem:L4} and Part~(B) of Theorem~\ref{thm:T1} are all conditional on it, which
  makes this the largest single conditionality the ledger carries.
\item \emph{Whether an equilibrium in which the blockholder chooses the fully separating plan exists
  on a given calibration.} Menus satisfying Assumption~\ref{asm:A7p} are fully separating on the
  flagged set, so the burden did not disappear when satisfiability was resolved; it moved to
  incentive compatibility, which Proposition~\ref{prop:P1} does not settle. Relatedly, that
  proposition does not claim that an equilibrium satisfying Assumption~\ref{asm:A8} exists --- only
  that Assumption~\ref{asm:A8} holding \emph{at} an exhibited equilibrium puts both cells on path.
\item \emph{Whether the continuity clause of Assumption~\ref{asm:A6} holds for the declared
  construction of $\Theta$.} This was answered in substance on 2026-08-27 and the locus was
  corrected --- see Remark~\ref{rem:A6record} --- and what remains open is scoped: (a) whether a
  constructive $\Theta$, or the $t$-constrained Kakutani route already on file, which removes the
  continuity half of the hypothesis, replaces that hypothesis in the statement, the repair being
  identified rather than executed; (b) the complementary menu class, in which every middle plan's
  histories are shared with a survivor --- the implemented menu is one such for its collapse face
  --- where the collapse-face clause may be satisfiable; and (c) nonexistence, which is neither
  claimed nor shown anywhere.
\end{enumerate}
